\section{The Six Birds dictionary for mathematical closures}
\label{sec:dictionary}

The Six Birds framework \cite{SixBirdsFramework} names six roles that recur whenever a detailed description is compressed into a coarser one that is still supposed to obey laws.
In this paper the roles are a naming convention: each one marks a piece of data that must be supplied explicitly when a discrete protocol is compressed into a stable mathematical object.
The dictionary is adopted, not derived.
For every exhibit we supply the carrier, the comparison, the packaging, and the audit by hand; none of them follows from the six labels, and the framework's own structural hypotheses (such as its bounded-interface assumption) are not tested by the experiments here.

\subsection{The six roles}
In mathematics the objects playing these roles are usually families indexed by a refinement parameter (meshes, truncations, partitions), error terms and residual norms, quotient or completion maps, and coherence conditions such as linearity, locality, invariance, or compatibility with products.

\begin{remark}[The six roles in mathematical form]\label{rem:six-primitives}
Let $\mathcal P$ be a set of protocol objects (discrete descriptions), refined along a stage parameter $j$ (or a step size $h\downarrow 0$).
\begin{enumerate}
  \item \Pfive\ (Packaging). An equivalence relation $\sim$ on $\mathcal P$ identifies descriptions that should count as the same object, with quotient map $\Pi:\mathcal P\to\mathcal P/\!\sim$.
  The map $\Pi$ changes type, so it is not itself an idempotent endomap.
  The idempotent object is the \emph{saturation} of sets of descriptions,
  \[
    \mathrm{cl}(A)\;:=\;\Pi^{-1}\bigl(\Pi(A)\bigr),\qquad A\subseteq\mathcal P,
  \]
  which is extensive, monotone, and idempotent, that is, a closure operator \cite{DaveyPriestley2002}.
  A finer equivalence identifies less, so its saturation is smaller and has more fixed points.
  In analytic instances, $\sim$ is ``equal up to a defect that vanishes under refinement'' (Section~\ref{subsec:packaging}); in algebraic instances it is a congruence.

  \item \Psix\ (Accounting). A ledger is a designated family of nonnegative diagnostics: error bounds, residual norms, route-mismatch statistics.
  In this paper the ledger records feasibility, that is, whether the compression stays coherent under refinement; it is not an arrow-of-time certificate.
  Refinement orders can support monotone information measures, but they do not by themselves make an approximation error decrease. Whether the ledger shrinks is something we test, never something we assume.

  \item \Pfour\ (Staging). The refinement parameter $j$ or $h$ is the stage. Staging is nontrivial when refinement genuinely increases resolution rather than merely reparametrizing.

  \item \Ptwo\ (Constraints). Coherence conditions imposed on the packaged layer: linearity, a product rule, an invariance, or compatibility with an algebraic operation.

  \item \Pone\ (Operator rewrite). A micro-operator $F:\mathcal P\to\mathcal P$ induces a packaged operator $F^\sharp([p]):=[F(p)]$ exactly when
  \[
    p\sim q\ \Longrightarrow\ F(p)\sim F(q).
  \]
  When this fails, one must change the comparison, rescale or otherwise rewrite $F$ (for example by the factor $h^{-1}$ that turns a difference into a difference quotient), or accept that no packaged operator exists.

  \item \Pthree\ (Protocol). When two admissible routes produce objects of the same type (``refine, then apply'' versus ``apply, then refine''), their discrepancy is a route-mismatch defect $\RM$. In stable regimes it should shrink under refinement.
\end{enumerate}
\end{remark}

\subsection{``A mathematics is A theory''}
The framework organizes an instantiation as a \emph{theory package}: a carrier, a lens, a completion rule, and an audit \cite{SixBirdsFramework}.
We keep the same shape and read the audit as a defect ledger.

\begin{definition}[Mathematics theory package]\label{def:math-theory-package}
A \emph{mathematics theory package} is a tuple $\mathcal T=(Z,\ f,\ \Sigma_f,\ E,\ \mathcal D)$ where
\begin{itemize}
  \item $Z$ is a carrier of discrete protocol objects at a chosen stage (grid functions, truncated sums or products, finite compositions, sequences of rationals);
  \item $f:Z\to X$ is a \emph{lens}, a coarse description map, and $\Sigma_f$ is the family of predicates on $Z$ that are visible through $f$, that is, constant on the fibers of $f$;
  \item $E:\mathcal V\to\mathcal V$ is a \emph{completion} or packaging endomap on a description space $\mathcal V$ associated with $Z$, idempotent or approximately idempotent in a stated norm, whose fixed points are the objects the package treats as complete;
  \item $\mathcal D$ is a \emph{defect ledger}, a designated family of nonnegative diagnostics together with the behaviour required of them under refinement (typically, decay to zero along the stage parameter). That behaviour is stated as a hypothesis or tested; it is not built into the definition.
\end{itemize}
\end{definition}

\begin{remark}[The title phrase]\label{rem:math-is-theory}
``A mathematics is A theory'' is meant in the following sense.
To introduce a stable mathematical layer is to specify (i) which discrete protocol objects are staged, (ii) which identifications are permitted under refinement, (iii) which completion or packaging map is used, and (iv) which defect ledger certifies that the packaged layer is coherent.
\end{remark}

\subsection{Notation}
We write $h>0$ for a step size and $h\downarrow 0$ for refinement; equivalently, a depth index $j\to\infty$.
Families of discrete objects carry stage subscripts ($A_h$, $F_h$, $z_h$).
Throughout, a route mismatch compares two routes $R_1(h)$ and $R_2(h)$ that produce outputs on a common comparison space:
\begin{equation}\label{eq:rm-generic}
  \RM(h)\;=\;\frac{\|R_1(h)-R_2(h)\|}{\|R_1(h)\|+\varepsilon_0},
\end{equation}
where $\varepsilon_0>0$ is a small constant that only guards against division by zero.
This is a \emph{relative} diagnostic.
Section~\ref{sec:exhibit-holonomy} instantiates it for a change of coordinates, Section~\ref{sec:exhibit-calculus} for two quadrature rules, and Section~\ref{sec:exhibit-prime} for two staged descriptions of~$\zeta$.
